\begingroup
\renewcommand{\arraystretch}{1.45}
\fancyhead[L]{\small ARITHMETIC–GEOMETRIC MEAN.}
\Needspace{28\baselineskip}
\AGMArticle{23}
[If the same algorithm as before for $\alpha,\allowbreak \beta,\allowbreak \gamma,\allowbreak \delta$ is now formed for $A,\allowbreak B,\allowbreak C,\allowbreak D$ and $A=\gamma,\allowbreak \ B=\delta$ is set, then evidently, for every $n$, $A^{\mathrm n}=\gamma^{\mathrm n},\allowbreak \ B^{\mathrm n}=\delta^{\mathrm n},\allowbreak \ C^{\mathrm n}=\alpha^{\mathrm n},\allowbreak \ D^{\mathrm n}=\beta^{\mathrm n}$, and hence for the limits $K,\allowbreak H$ corresponding to
\SourceLeaf{405}
$x,\allowbreak \eta$: $\frac{x}k=\frac Kk\cdot y^{\frac12}\cdot H,\allowbreak \ \frac{x}k y^{\frac12}\eta=\frac Kk$; the functional equations are therefore
\Needspace{14\baselineskip}
\[
\begin{aligned}\dfrac{R(y,\dfrac1{y\eta})}P
&=\dfrac{S(y,\dfrac1{y\eta})}Q
\\[7pt]&=\dfrac{P(y,\dfrac1{y\eta})}R
\\[7pt]&=\dfrac{Q(y,\dfrac1{y\eta})}S
\\[7pt]&=y^{-\frac14}\eta^{-\frac12}\end{aligned}
\]

If, however, $A=\beta,\allowbreak \ B=\alpha$ is set, then $C=-\delta,\allowbreak \ D=-\gamma$, and for every $n$ equal to or greater than unity: $A^{\mathrm n}=\alpha^{\mathrm n},\allowbreak \ B^{\mathrm n}=\beta^{\mathrm n},\allowbreak \ C^{\mathrm n}=\gamma^{\mathrm n},\allowbreak \ D=\delta^{\mathrm n}$; thus one obtains:
\Needspace{14\baselineskip}
\[
\begin{aligned}\dfrac{Q(y,-\eta)}P
&=\dfrac{P(y,-\eta)}Q
\\[7pt]&=-\dfrac{iS(y,-\eta)}R
\\[7pt]&=-\dfrac{iR(y,-\eta)}S\\[7pt]&=1\end{aligned}
\]

\Needspace{22\baselineskip}
Finally, if $2\surd {A''}=\surd \alpha+\surd \beta,\allowbreak \ 2\surd {C''}=\surd \alpha-\surd \beta$ is set, then for every $n\geqq1$
\[
\begin{gathered}
2A^{\mathrm n+1}=\surd {a^{\mathrm n}\alpha^{\mathrm n}}+\surd {b^{\mathrm n}\beta^{\mathrm n}},\\[7pt]
2C^{\mathrm n+1}=\surd {a^{\mathrm n}\alpha^{\mathrm n}}-\surd {b^{\mathrm n}\beta^{\mathrm n}}\\[7pt]
2B^{\mathrm n+1}=\surd {b^{\mathrm n}\alpha^{\mathrm n}}+\surd {a^{\mathrm n}\beta^{\mathrm n}},\\[7pt]
2D^{\mathrm n+1}=\surd {b^{\mathrm n}\alpha^{\mathrm n}}-\surd {a^{\mathrm n}\beta^{\mathrm n}}\\[7pt]
2\surd {A^{\mathrm n+2}}=\surd {\alpha^{\mathrm n}}+\surd {\beta^{\mathrm n}},\\[7pt]
2\surd {C^{\mathrm n+2}}=\surd {\alpha^{\mathrm n}}-\surd {\beta^{\mathrm n}},\\[7pt]
4A^{\mathrm n+1}C^{\mathrm n+1}=c^{\mathrm n}\gamma^{\mathrm n}\\[7pt]
\left(\dfrac Kk\right)^4=\dfrac{x}k,\\[7pt] HH=\eta
\end{gathered}
\]
and hence:
\Needspace{12\baselineskip}
\[
\begin{aligned}
2P(y^4,\eta\eta)&=P+Q,\\[6pt]2R(y^4,\eta\eta)&=P-Q\\[6pt]
2P(yy,\eta)^2&=pP+qQ,\\[6pt]2R(yy,\eta)^2&=pP-qQ\\[6pt]
2Q(yy,\eta)^2&=qP+pQ.\\[6pt]2S(yy,\eta)^2&=qP-pQ
\end{aligned}
\]
\Needspace{12\baselineskip}
\[
\begin{aligned}
p(\surd  y)\cdot P(\surd  y,\eta)&=PP+RR,\\[6pt]
q(\surd  y)\cdot P(\surd  y,\eta)&=QQ-SS,\\[6pt]
r(\surd  y)\cdot R(\surd  y,\eta&=2PR\\[6pt]
q(\surd  y)\cdot Q(\surd  y,\eta)&=PP-RR,\\[6pt]
p(\surd  y)\cdot Q(\surd  y,\eta)&=QQ+SS,\\[6pt]
r(\surd  y)\cdot S(\surd  y,\eta&=2QS
\end{aligned}
\]
where once again those functions to which no arguments are appended refer to the simple $y$ and $\eta$.

The new algorithm thus obtained, in which one passes from the functions $P,\allowbreak Q,\allowbreak R,\allowbreak S$ with the arguments $y,\allowbreak \eta$ to the functions with the arguments $yy,\allowbreak \eta$, is evidently the one used by \textsc{Gauss} in Article 16 of the \textit{Determinatio attractionis quam in punctum quodvis positionis datae exerceret planeta etc.}

The relations between the functions referring to the two arbitrary values $\xi,\allowbreak \eta$ of the second argument and those functions with the second arguments $\xi\eta$ and $\frac\xi\eta$ can be derived in various ways with
\SourceLeaf{406}
the aid of the algorithm established for $\alpha,\allowbreak \beta$. The method in which the developments are least lengthy seems to be the one referring to functions in which, alongside the second arguments $\xi\eta$ and $\frac{\xi}{\eta}$, the first argument is the square of the first argument of the functions that have $\xi,\allowbreak \eta$ as their second argument.]

\Needspace{22\baselineskip}
\AGMArticle{24}
[In the algorithm of the arithmetic–geometric mean, the terms arising through backward extension yield, with the aid of the equations
\Needspace{12\baselineskip}
\[
\begin{aligned}\dfrac{a_{\mathrm n}}{{}^{\mathrm n}\!a}&=\dfrac{c_{\mathrm n}}{{}^{\mathrm n}\!c}\\[7pt]&=\dfrac{b_{\mathrm n}}{{}^{\mathrm n}\!b}\\[7pt]&=\dfrac{1}{2^{\mathrm n}}\end{aligned}
\]
\Needspace{28\baselineskip}
the terms $a_{\mathrm n},\allowbreak c_{\mathrm n},\allowbreak b_{\mathrm n}$, which depend on $a,\allowbreak c$ in the same way as $a^{\mathrm n},\allowbreak b^{\mathrm n},\allowbreak c^{\mathrm n}$ depend on $a,\allowbreak b$. If, accordingly, in the combined algorithm $b$ is interchanged with $c$ and $\beta$ with $\gamma$ at the same time, while $\alpha$ is left unchanged, then $\delta$ changes only its sign; we therefore obtain an algorithm that depends on $a,\allowbreak c,\allowbreak \alpha,\allowbreak \gamma$ in the same way as the one considered so far depends on $a,\allowbreak b,\allowbreak \alpha,\allowbreak \beta$, if we set
\[
\begin{gathered}
\alpha_0=\alpha,\\[7pt]\gamma_0=\beta,\\[7pt]\beta_0=\gamma,\\[7pt]\delta_0=-\delta\\[8pt]
\begin{aligned}
\alpha_1&=\dfrac{1}{a_1}\left(\dfrac{\alpha_0+\gamma_0}{2}\right)^2
\\[6pt]&=\dfrac{1}{b_1}\left(\dfrac{\beta_0-\delta_0}{2}\right)^2,
\\[7pt]\gamma_1&=\dfrac{1}{c_1}\alpha_0\gamma_0\\[8pt]
\beta_1&=\dfrac{1}{b_1}\left(\dfrac{\alpha_0-\gamma_0}{2}\right)^2
\\[6pt]&=\dfrac{1}{a_1}\left(\dfrac{\beta_0+\delta_0}{2}\right)^2,
\\[7pt]\delta_1&=\dfrac{1}{c_1}\beta_0\delta_0
\end{aligned}
\end{gathered}
\]
\Needspace{20\baselineskip}
and so on up to the limits:
\[
\begin{gathered}
\lim a_{\mathrm n}=\lim c_{\mathrm n}=\mathrm{M}(a,c)=l\\[7pt]
\lim{}^{2^{\mathrm n}}\!\surd {\dfrac{b_{\mathrm n}}{4l}}=e^{-\dfrac{\pi}{2}\dfrac{\mathrm{M}(a,c)}{\mathrm{M}(a,b)}}=\surd  z\\[7pt]
\lim{}^{2^{\mathrm n}}\!\surd {\dfrac{\alpha_{\mathrm n}}{l}}=\lim{}^{2^{\mathrm n}}\!\surd {\dfrac{\gamma_{\mathrm n}}{l}}=\dfrac{\lambda}{l}\\[7pt]
{}^{2^{n-1}}\!\surd {\left[\surd {\dfrac{\beta_{\mathrm n}}{4l}}\mathbin{\pm}\surd {\dfrac{\delta_{\mathrm n}}{4l}}\right]}=\dfrac{\lambda}{l}z^{\frac12}\zeta^{\pm1}
\end{gathered}
\]
Thus $\alpha,\allowbreak \gamma,\allowbreak \beta,\allowbreak -\delta$ are functions of $l,\allowbreak z,\allowbreak \lambda,\allowbreak \zeta$ of the same kind as $\alpha,\allowbreak \beta,\allowbreak \gamma,\allowbreak \delta$ are of $k,\allowbreak y,\allowbreak x,\allowbreak \eta$; and since $\frac{l}{\lambda}$ depends only on $y$ and $\eta$, in precisely the same way as $\frac{\lambda}{l}$ depends on $\eta$ and $y$, for $\frac{\log y}{\pi}=\frac{\pi}{\log z}$,

\SourceLeaf{407}

\Needspace{22\baselineskip}
\[
\begin{aligned}\dfrac{P(z,\zeta)}{P(y,\eta)}&=\dfrac{R(z,\zeta)}{Q(y,\eta)}\\[7pt]&=\dfrac{Q(z,\zeta)}{R(y,\eta)}\\[7pt]&=\dfrac{iS(z,\zeta)}{S(y,\eta)}\\[7pt]&=T(y,\eta)\\[7pt]&=\dfrac{1}{T(z,\zeta)}\end{aligned}
\]
\Needspace{10\baselineskip}
must hold. Here, besides $T$, $\zeta$ must still be determined as a function of $y,\allowbreak \eta$, and specifically $\zeta$ so that $y,\allowbreak \eta,\allowbreak \zeta$ can be interchanged respectively with $z,\allowbreak \zeta,\allowbreak \eta$, thus so that
\[
\begin{aligned}\dfrac{\log y}{\pi}&=\dfrac{\pi}{\log z}\\[7pt]&=\varphi[(\log\eta)^2,(\log\zeta)^2]\\[7pt]&=\dfrac{1}{\varphi[(\log\zeta)^2,(\log\eta)^2]}\end{aligned}
\]
holds.

\Needspace{24\baselineskip}
Here the same signs have been taken for $P,\allowbreak Q,\allowbreak R$ as square roots of $\frac{\alpha}{x},\allowbreak \frac{\beta}{x},\allowbreak \frac{\gamma}{x}$, because by Article 18 $\eta$ and $\zeta$ take the value $1$ at the same time, and hence by Articles 22 and 17
\[
\begin{aligned}\dfrac{P(z,1)}{P(y,1)}&=\dfrac{R(z,1)}{Q(y,1)}\\[7pt]&=\dfrac{Q(z,1)}{R(y,1)}\\[7pt]&=T(y,1)\\[7pt]&=\dfrac{1}{T(z,1)}\\[7pt]&=\dfrac{pz}{py}\\[7pt]&=\dfrac{rz}{qy}\\[7pt]&=\dfrac{qz}{ry}\\[7pt]&=\surd {\dfrac{-\log y}{\pi}}\\[7pt]&=\surd {\dfrac{\pi}{-\log z}}\end{aligned}
\]
holds; the sign of the square root used here is the one for which the real part is positive.

If the second arguments $\zeta,\allowbreak \eta$ are changed into their reciprocals, then $P,\allowbreak Q,\allowbreak R$ remain unchanged and $S$ changes only its sign; thus
\Needspace{5\baselineskip}
\[
T(y,\eta)=T(y,\tfrac{1}{\eta})
\]
and a corresponding condition holds for $y$ as a function of $\eta,\allowbreak \zeta$, as is intended to be indicated by the form of the $\varphi$ function given above.

\Needspace{18\baselineskip}
If, in the equation by which $T$ was introduced, $P,\allowbreak Q,\allowbreak R,\allowbreak S$ are transformed so that the arguments $z,\allowbreak \zeta,\allowbreak y,\allowbreak \eta$ pass first into $z,\allowbreak -\zeta,\allowbreak y,\allowbreak \frac{1}{y\eta}$ and then into $\surd  z,\allowbreak \zeta,\allowbreak yy,\allowbreak \eta\eta$, and account is taken of the fact that
\[
\begin{aligned}\dfrac{p\surd  z}{pyy}&=\dfrac{r\surd  z}{qyy}\\[7pt]&=\dfrac{q\surd  z}{ryy}\\[7pt]&=\surd {\dfrac{-\log yy}{\pi}}\\[7pt]&=\surd {\dfrac{\pi}{-\log\surd  z}}\end{aligned}
\]
holds, one obtains
\Needspace{24\baselineskip}
\[
\begin{gathered}\begin{aligned}\dfrac{Q(z,-\zeta)}{R(y,\dfrac{1}{y\eta})}&=\dfrac{iS(z,-\zeta)}{S(y,\dfrac{1}{y\eta})}\\[7pt]&=\dfrac{P(z,-\zeta)}{R(y,\dfrac{1}{y\eta})}\\[7pt]&=\dfrac{R(z,-\zeta)}{Q(y,\dfrac{1}{y\eta})}\\[7pt]&=y^{\frac14}\eta^{\frac12}T(y,\eta)\\[7pt]&=T(y,\tfrac{1}{y\eta})\end{aligned}\\[10pt]\begin{aligned}\dfrac{P(\surd  z,\zeta)}{P(yy,\eta\eta)}&=\dfrac{R(\surd  z,\zeta)}{Q(yy,\eta\eta)}\\[7pt]&=\dfrac{Q(\surd  z,\zeta)}{R(yy,\eta\eta)}\\[7pt]&=\dfrac{iS(\surd  z,\zeta)}{S(yy,\eta\eta)}\\[7pt]&=2\surd {\dfrac{\pi}{-\log yy}}\cdot T(y,\eta)^2\\[7pt]&=T(yy,\eta\eta)\end{aligned}\end{gathered}
\]

\SourceLeaf{408}

\Needspace{22\baselineskip}
\[
\begin{gathered}\begin{aligned}\dfrac{\log y}{\pi}&=\dfrac{\pi}{\log z}\\[7pt]&=\varphi[(\log\eta)^2,(\log\zeta)^2]\\[7pt]&=\varphi[(\log\eta)^2,(\pi i+\log\zeta)^2]\end{aligned}\\[10pt]\begin{aligned}\dfrac{\log yy}{\pi}&=\dfrac{\pi}{\log\surd  z}\\[7pt]&=2\varphi[(\log\eta)^2,\log\zeta)^2]\\[7pt]&=\varphi[(\log\eta\eta)^2,(\log\zeta)^2]\end{aligned}\end{gathered}
\]
and hence
\Needspace{8\baselineskip}
\[
\begin{gathered}
T(y,\eta)=\surd {\dfrac{-\log y}{\pi}}\cdot e^{\frac{(\log\eta)^2}{4\log y}}\\[7pt]
\dfrac{\log y}{\pi}=\dfrac{\pi}{\log z}=\mathbin{\pm}i\cdot\dfrac{\log\eta}{\log\zeta}
\end{gathered}
\]
Here the sign of the square root is to be taken so that its real part is positive, but the sign of $\mathbin{\pm}i\frac{\log\eta}{\log\zeta}$ so that the real part of this expression is negative.]

\clearpage
\endgroup
